%%%%%%%%%%%%%%%%%%%%%%%%%%%%   SECTION 0 - LICENSE %%%%%%%%%%%%%%%%%%%%%%%%%%%%%%%%%%%%%%%%%
% Copyright 2018 Terrence P. Murphy and Rosalie D'Amico.
% This file may be distributed and/or modified under the conditions of the LaTeX Project 
% Public License (“LPPL”), either version 1.3c of this license or (at your option) any 
% later version. The latest version of this license is at:
%     http://www.latex-project.org/lppl.txt.
% This file is author-maintained and  is one of the files that, collectively, make up the 
% MAKECOOKBOOK bundle available at https://ctan.org/pkg/makecookbook
% For additional information, please see the associated README file.
%%%%%%%%%%%%%%%%%%%%%%%%%%%%   SECTION 1 - INITIAL SETUP %%%%%%%%%%%%%%%%%%%%%%%%%%%%%%%%%%%
\raggedbottom				% Q320739 and Q65355
\emergencystretch 3em     	% can also use for only parts of a document
%\tracingmacros=1
%\tracingcommands=1 % or =2

%%%%%%%%%%%%%%%%%%%%%%%%%%%%%%%%%%%%%%%%%%%%%%%%%%%%%%%%%%%%%%%%%%%%%%%%%%%%%%%%%%%%%%%%%%%
% We use the TeX \newif macro to choose betweeen the makecookbook code and the code for
% our own cookbook.  One of the key differences is our use of commercial fonts (Adobe Garamond 
% Pro and Adobe Myriad Pro) vs. makecookbook's free fonts (EB Garamond and Lato).
%%%%%%%%%%%%%%%%%%%%%%%%%%%%%%%%%%%%%%%%%%%%%%%%%%%%%%%%%%%%%%%%%%%%%%%%%%%%%%%%%%%%%%%%%%%
\newif\ifMakeCookbook    \MakeCookbooktrue
% \newif to indicate whether \geometry (and other) settings are for draft or final form
\newif\ifCookbookDraft   \CookbookDrafttrue
% newif to indicate whether we are using the hyperref/bookmark packages
\newif\ifHyperRef        \HyperReftrue

%************************* FIRST PACKAGES **************************************************
% We load the packages that are used for multiple purposes.  The single-purpose packages
% will be loaded with their associated code to make their purpose clearer.
%*******************************************************************************************
\usepackage{%
	xparse, 			% For \NewDocumentCommand and related
	enumitem,			% improvements to the standard list environments
	pgfkeys,			% For Key=Value arguments to commands
	graphicx}           % for including photos = \includegraphics

\ifCookbookDraft
	\usepackage[base]{babel}
	\usepackage{lipsum, layout} %
\fi
%\usepackage{indentfirst}   % use if you want first paragraphs indented; see Q 39227 (Warner)

%%%%%%%%%%%%%%%%%%%%%%%%%%%%   SECTION 2 - FONTS %%%%%%%%%%%%%%%%%%%%%%%%%%%%%%%%%%%%%%%%%%%

%************ SECTION 2A - FONT PACKAGES ***************************************************
\usepackage{fontspec}           % Needed for our font handling -- how about [no-math] ??
\usepackage{microtype}          % Highly advised-READ UP to understand better; see Q 66052 and S.K. link

%************ SECTION 2B - ESTABLISH ALL FONTS TO BE USED **********************************
\ifMakeCookbook
	\setmainfont{ebgaramond}[   			% EB Garamond
		Extension=.ttf,
		UprightFont=*-regular,                  % this face includes small caps
		ItalicFont=*-italic,
		BoldFont=*-bold,
		BoldItalicFont=*-bolditalic,
		FontFace={sb}{n}{*-semibold},           % this face includes small caps   
		FontFace={sb}{it}{*-semibolditalic},
		Ligatures=TeX,
		Numbers=Lining]				 			% !!! consider Numbers=OldStyle
\else
	\setmainfont{agaramondpro}[   			% Abobe Garamond Pro 
		Extension=.otf,
		UprightFont=*-regular,                  % this face includes small caps
		ItalicFont=*-italic,
		BoldFont=*-bold,
		BoldItalicFont=*-bolditalic,
		FontFace={sb}{n}{*-semibold},           % this face includes small caps   
		FontFace={sb}{it}{*-semibolditalic},
		Ligatures=TeX,
		Numbers=Lining]
\fi

\ifMakeCookbook
	\setsansfont{lato}[     			% Lato
		Extension=.ttf,
		UprightFont=*-regular,       		% \mdseries
		ItalicFont=*-italic,
		BoldFont=*-bold,             		% \bfseries
		BoldItalicFont=*-bolditalic,
		FontFace={sb}{n}{*-semibold},
		FontFace={sb}{it}{*-semibolditalic},
		FontFace={k}{n}{*-black},
		FontFace={k}{it}{*-blackitalic},
		Ligatures=TeX,
		Numbers=Lining]
\else
	\setsansfont{myriadpro}[     		% Adobe Myriad Pro
		Extension=.otf,              		% None of the font faces include small caps
		UprightFont=*-regular,       		% \mdseries
		ItalicFont=*-it,
		BoldFont=*-bold,             		% \bfseries
		BoldItalicFont=*-boldit,
		FontFace={sb}{n}{*-semibold},
		FontFace={sb}{it}{*-semiboldit},
		FontFace={k}{n}{*-black},
		FontFace={k}{it}{*-blackit},
		Ligatures=TeX,
		Numbers=Lining]						% the OldStyle '1' is not too good for ingredients
\fi

\newfontfamily\MyScriptFont{Italianno}[	% Italianno script font
	Extension=.ttf,
	UprightFont=*-Regular,                  % This is the only face for this font
	Ligatures=TeX]                          % Otherwise, don't get -- converted to en dash, etc. Q 53413
% (and \textendash doesn't work that well with Italianno)

%************ SECTION 2C - HANDLE THE SPECIAL FONT FACES ***********************************
\NewDocumentCommand \sbseries {}{\fontseries{sb}\selectfont} % egreg uses \DeclareRobustCommand
\DeclareTextFontCommand{\textsb}{\sbseries}
\NewDocumentCommand \kseries {}{\fontseries{k}\selectfont}   % egreg uses \DeclareRobustCommand
\DeclareTextFontCommand{\textk}{\kseries}

%************ SECTION 2D - SET COMMANDS TO SELECT FONTS ************************************
\NewDocumentCommand \FontSteps          {}{\rmfamily\mdseries}
\NewDocumentCommand \FontStepsDefault   {}{\rmfamily\mdseries}
\NewDocumentCommand \FontIngredients    {}{\sffamily\sbseries}
\NewDocumentCommand \FontIngDefault  	{}{\sffamily\sbseries}
\NewDocumentCommand \FontIngTitle		{}{\sffamily\bfseries\itshape}  % consider \kseries
\NewDocumentCommand \FontChapterIntro   {}{\rmfamily\fontsize{14}{16.8}\selectfont}
\NewDocumentCommand \FontRecipeStory    {}{\rmfamily\fontsize{13}{15.6}\selectfont}
\NewDocumentCommand \FontCopyrightPage  {}{\rmfamily\fontsize{9}{11}\selectfont}
\NewDocumentCommand \FontChapterLabel   {}{\rmfamily\fontsize{14.4}{18}\selectfont}
\NewDocumentCommand \FontChapterTitle   {}{\MyScriptFont\fontsize{40}{48}\selectfont}
\NewDocumentCommand \FontRecipeName     {}{\MyScriptFont\fontsize{24}{29}\selectfont}
\NewDocumentCommand \FontLettrineText   {}{\rmfamily\scshape}	% consider adding \sbseries
\NewDocumentCommand \FontTitleColorBox  {}{\rmfamily\Large\bfseries}
\NewDocumentCommand \FontTitlepageTitle {}{\MyScriptFont\fontsize{40}{48}\selectfont}
\NewDocumentCommand \FontTitlepageAuthor{}{\rmfamily\sbseries\scshape\fontsize{14.4}{18}\selectfont}

%************ SECTION 2E - SPECIAL HANDLING OF FRACTIONS ***********************************
\def\fr#1/#2 {{\addfontfeatures{Fractions=On}#1/#2} }	%
\def\frx#1/#2 {{\addfontfeatures{Fractions=On}#1/#2}}	%

\usepackage{xfrac}

%************ SECTION 2F - SPECIAL GLYPHS NEEDED IN COOKBOOK *******************************
\NewDocumentCommand \TextDegree	{}{{\char176}}		% or ^^^^00B0
\NewDocumentCommand \Copyright	{}{{\char169}}		% or ^^^^00A9
\NewDocumentCommand \TextBullet	{}{{\char8226}} 	% or ^^^^2022
\NewDocumentCommand \CtrDot		{}{{\char183}}		% or ^^^^00B7
% The following replaces $\cdots$ -- it is essentially identical and uses our fonts:
\NewDocumentCommand \CtrDots{}{{\CtrDot\kern 0.2em\CtrDot\kern 0.2em\CtrDot\kern 0.2em}}%

% \ChefHat and \Oven are not actually glyphs - they were created (further below) 
% with \MakeChefHat and \MakeOven using TIKZ
\newsavebox{\HatBox}
\AtBeginDocument{\savebox{\HatBox}[\hatwidth]{\MakeChefHat}}%
\NewDocumentCommand \ChefHat	{}{\usebox{\HatBox}}%
\NewDocumentCommand \ChefNote   {}{{\raisebox{.4ex}{\ChefHat}}}

\newsavebox{\OvenBox}
\AtBeginDocument{\savebox{\OvenBox}[\ovenwidth]{\MakeOven}}%
\NewDocumentCommand \Oven		{}{\usebox{\OvenBox}}%	

%************ SECTION 2G - FONT SIZE INFORMATION NEEDED IN CODE FURTHER BELOW ***************
\def\StdIFontSize{10.95}		% Needed in IngredientsAndSteps environment
\def\StdSFontSize{10.95}		% Needed in IngredientsAndSteps environment
\def\StdIBaseline{13.6}			% Needed in IngredientsAndSteps environment
\def\StdSBaseline{13.6}			% Needed in IngredientsAndSteps environment
\def\RStoryFontBaseline{15.6pt} % Needed in \RecipeStory command; value from \FontRecipeStory
\def\CIntroFontBaseline{16.8pt} % Needed in \ChapterIntro command; value from \FontChapterIntro

%%%%%%%%%%%%%%%%%%%%%%%%%%%%   SECTION 3 - PACKAGES %%%%%%%%%%%%%%%%%%%%%%%%%%%%%%%%%%%%%%%%         
%************ SECTION 3A - GEOMETRY ********************************************************
\usepackage{geometry} %

\ifCookbookDraft
	\geometry{paper=letterpaper,      % the physical paper size during draft mode
		layoutsize={8in,10in},            % always use intended final paper size for layout	
		layouthoffset=0.25in,             % center the "layout" horizontally
		layoutvoffset=0.5in,              % center the "layout" vertically
		%showframe,                       % use when needed
		showcrop}
\else
	\geometry{papersize={8in,10in}}   % the physical paper size in final production mode
\fi
% The following are the same in draft and final:
\geometry{nomarginpar,            % do not reserve space for margin notes
	bindingoffset=0.375in,
	inner=0.5in,
	outer=1in,
	top=0.75in,
	bottom=1in,
	footskip=40pt}                    % default seems to be 27pt

%************ SECTION 3B - IMAKEIDX PACKAGE ************************************************
\usepackage{imakeidx}				% supports creation of an index (here, a recipe index)

%%% COMMENT OUT THE FOLLOWING CODE IF YOU WANT THE STANDARD BOOK CLASS INDEX FORMAT %%%%%%
\makeatletter
\def\@idxitem{\par\hangindent 10pt} % NOT needed unless you want to fine tune hanging indent
\newcommand{\betterdotfill}%
{\leavevmode \nobreak\cleaders \hb@xt@ .44em{\hss .\hss }\hskip .5em plus 1fill \kern \z@}
\makeatother
\makeindex[options=-s cb-idxstyle]		% use cb-idxstyle.ist for style; Q 132465 & Q 396898
%%% END OF SPECIAL INDEX FORMAT %%%%%%%%%%%%%%%%%%%%%%%%%%%%%%%%%%%%%%%%%%%%%%%%%%%%%%%%%%

%************ SECTION 3C - LETTRINE PACKAGE ************************************************
\usepackage{lettrine}			% Sets first character to multi-line height with various options
\renewcommand{\LettrineFontHook}{\MyScriptFont\color{clrLettrineBig}}		% for the BIG first letter
\renewcommand{\LettrineTextFont}{\color{clrLettrineSmall}\FontLettrineText}	% for the following text
\renewcommand{\DefaultOptionsFile}{cb-lettrine.cfl}						% for the lettrine package

%************ SECTION 3D - FANCYHDR PACKAGE ************************************************
\usepackage{emptypage}	% per Q 360739 replaces \clearpage{\pagestyle{empty}\cleardoublepage} which
% would have to be issued at the end of each chapter and possibly the last 
% page of the frontmatter and as needed in the backmatter
\usepackage{fancyhdr}

\fancypagestyle{plain}{ 				 % Need to redefine the "plain" style
	\fancyhf{}                             % clear the header and footer
	\renewcommand{\headrulewidth}{0pt}     % use 0 to disable header ruler line
	\renewcommand{\footrulewidth}{0.2pt}
	\fancyfoot[RO, LE] {Page \thepage}
}

\makeatletter % must be OUTSIDE of the command - see Q 444532
\fancypagestyle{main}{  % identical to plain, except non-\chapter pages in mainmatter                 
	\fancyhf{}
	\renewcommand{\headrulewidth}{0pt}
	\renewcommand{\footrulewidth}{0.2pt}
	\fancyfoot[RO, LE] {Page \thepage}
	\fancyfoot[CE,CO]{\if@mainmatter \leftmark\fi}  % See Q340125
}
\makeatother

%************ SECTION 3E - TITLESEC PACKAGE ************************************************
\usepackage{titlesec}
\titleformat{\chapter}[display]
{\filleft\FontChapterLabel}{\chaptertitlename\ \thechapter}{1pt}
{\titlerule\vspace{1ex}\raggedleft\FontChapterTitle}
\titlespacing*{\chapter}{0pt}{20pt}{20pt}

%%%%%%%%%%%%%%%%%%%%%%%%%%%%   SECTION 4 - CERTAIN SIMPLE MACROS %%%%%%%%%%%%%%%%%%%%%%%%%%%

%************ SECTION 4A - COLORS **********************************************************
\usepackage{xcolor}
\definecolor{clrWhite}{cmyk}{0.00, 0.00, 0.00, 0.00} 	% true white
\definecolor{clrBackTip}{rgb}{1.0, 0.95, 0.95}			% red!5!white
\definecolor{clrFrameTip}{rgb}{0.75, 0.0, 0.0}			% red!75!black
\definecolor{clrBackCheffy}{rgb}{1.0, 1.0, 1.0}			% white
\definecolor{clrFrameCheffy}{rgb}{0.0, 0.0, 0.75}		% blue!75!black
\definecolor{clrBackNotes}{rgb}{1.0, 1.0, 1.0}			% white
\definecolor{clrFrameNotes}{rgb}{0.0, 0.75, 0.0}		% green!75!black
\definecolor{clrLettrineBig}{gray}{0.35}				% 0.0 = black; 1.0 = white
\definecolor{clrLettrineSmall}{gray}{0.35}				% 0.0 = black; 1.0 = white
\definecolor{clrIngTitle}{cmyk}{0.00, 1.00, 1.00, 0.00} % true red
\definecolor{clrEditNote}{cmyk}{0.00, 1.00, 1.00, 0.00} % true red
\definecolor{clrHyperRef}{gray}{0.50} % true red

%************ SECTION 4B - A FIX FOR PGFKEYS ***********************************************
\def\nl{\par}		% see Q 96247 for why we defined \nl for use with pgfkeys

%************ SECTION 4C - VARIOUS SIMPLE BUT USEFUL COMMANDS ******************************
\NewDocumentCommand \PreheatC{m}{Preheat oven to #1\Degrees convection.\thinspace\Oven}%
\NewDocumentCommand \PreheatR{m}{Preheat oven to #1\Degrees regular oven.\thinspace\Oven}%

\NewDocumentCommand \Tbl{o}{\IfNoValueTF{#1}{tablespoon }{tablespoon#1}}%
\NewDocumentCommand \tsp{o}{\IfNoValueTF{#1}{teaspoon }{teaspoon#1}}%
\NewDocumentCommand \Pd{o}{\IfNoValueTF{#1}{pound }{pound#1}}%
\NewDocumentCommand \Ounce{o}{\IfNoValueTF{#1}{ounce }{ounce#1}}%
\NewDocumentCommand \Degrees{o}{\IfNoValueTF{#1}{\TextDegree{} }{\TextDegree{}#1}}%

\NewDocumentCommand \AxB{m m o}{{#1\thinspace{x}\thinspace#2}\IfValueT{#3}{#3}}%
\NewDocumentCommand \AxBxC{m m m o}{%
	{#1\thinspace{x}\thinspace#2\thinspace{x}\thinspace#3}\IfValueT{#4}{#4}}%

\NewDocumentCommand \Inch{m}{{#1-inch}}%
\NewDocumentCommand \EditNote{m}{{{\color{clrEditNote} #1}}}%
\NewDocumentCommand \Quote{m}{{``#1''}}%

\NewDocumentCommand \IngredientsSeparator{}{{\FontStepsDefault \CtrDots\CtrDots}}%
\NewDocumentCommand \IngredientsSeparatorClear{}{{\vskip 8pt}}%
\NewDocumentCommand \SeparateParagraphs{}{{\vskip 5pt}}%
\NewDocumentCommand \Recipe{+m}{{\textit{\textsb{#1}}}}%

%************ SECTION 4D - BAKE UNTIL ******************************************************
\pgfkeys{
	/BakeUntil/.is family, /BakeUntil,
	default/.style = {Min = 0, Max = 0, TPick = 0, GBrown = 0, Until = \empty},
	Min/.estore in 		= \BkMin,
	Max/.estore in 		= \BkMax,
	TPick/.estore in 	= \BkTPick,
	GBrown/.estore in 	= \BkGBrown,
	Until/.estore in 	= \BkUntil,
}

\NewDocumentCommand \BakeUntil{m}{%
\pgfkeys{/BakeUntil, default, #1}%	
\newcount\MyCount		\MyCount = \BkMax%
\newcount\MyCase%
{Bake about \BkMin}%
\ifnum\MyCount > 0 {-\BkMax}\fi
{ minutes}%
%
\ifx\BkUntil\empty  \MyCase = 0 %
\else \MyCase = 1 \fi%
\ifnum\BkTPick = 1  \MyCase = 2 \fi%
\ifnum\BkGBrown = 1 \MyCase = 3 \fi%
% 
\ifcase\MyCase% 						'0' case - do nothing
\or {, \BkUntil}% 						'1' case - use the "Until" text
\or {, until toothpick tests clean}% 	'2' case - TPick
\or {, until golden brown}% 			'3' case - GBrown
\fi
}

%************ SECTION 4E - OTHER MISCELLANEOUS COMMANDS ************************************
\NewDocumentEnvironment{CopyrightPage}{} % See 31186
{\FontCopyrightPage\setlength{\parindent}{0pt}\par\vspace*{\fill}}
{\clearpage}

\NewDocumentCommand \Attribution{m}{{\hspace*\fill}{\textit{#1}}}%
\NewDocumentCommand \ThreeLines{}{%
	Lorem ipsum dolor sit amet, consectetuer adipiscing elit. Ut purus elit, vestibulum ut,
	placerat ac, adipiscing vitae, felis. Curabitur dictum gravida mauris. Nam arcu libero,
	nonummy eget, consectetuer id, vulputate a, magna. Donec vehicula augue eu neque.
}

\newcount\LineCount
\NewDocumentCommand \InsertHiddenLines{m}{%
	\LineCount = #1
		{%
			\newcount\foo \foo=0
			\loop				             % \loop starts the construct ended by \repeat
			\phantom{.}\par \advance \foo by 1
			\ifnum \foo < \LineCount \repeat % \repeat also “serves as” the \fi to the \ifnum
		}%
}

%************ SECTION 4F - CHAPTER INTRO ***************************************************
\newbox\IntroBox
\newbox\IntroBoxA
\NewDocumentCommand \ChapterIntro {o +m}  % arg1 = optional = lines for first part
{                                         % arg2 = story text
	\setbox\IntroBox\vbox
	{
		\FontChapterIntro
		{#2\par}      % the \par needed by \lettrine
	}
	\IfNoValueTF {#1}    % NoValue=TRUE or NoValue=FALSE?
	{\unvbox\IntroBox} % NoValue=TRUE, so flush the full IntroBox
	{                  % handle the NoValue=FALSE case
		\ifnum #1 > 0      % > 0 means we split the story into two part; print the first part here
				{                % = 0 means ALL of story is deferred
					\setbox\IntroBoxA=\vsplit\IntroBox to #1\dimexpr \CIntroFontBaseline % 
					\unvbox\IntroBoxA  % 
				}
		\fi
	}
}

\NewDocumentCommand \FinishChapterIntro{}%
{\ifvoid\IntroBox \else \smallskip\unvbox\IntroBox\par\medskip\fi} %

%%%%%%%%%%%%%%%%%%%%%%%%%%%%   SECTION 5 - RECIPE ENVIRONMENTS AND COMMANDS %%%%%%%%%%%%%%%%

%************ SECTION 5A - RECIPE NAME AND YIELD *******************************************
\pgfkeys{
	/RecipeNameAndYield/.is family, /RecipeNameAndYield,
	default/.style = {Name = 0, NoIdxName = 0, XRefLabel = \empty, Yield = \empty, %
			IndexA = \empty, IndexB = \empty, IndexC = \empty},
	Name/.estore in  	 = \RecipeName,
	NoIdxName/.estore in = \NoIndexName,
	IndexA/.estore in    = \IdxA,
	IndexB/.estore in    = \IdxB,
	IndexC/.estore in    = \IdxC,
	XRefLabel/.estore in = \XRefLbl,
	Yield/.estore in 	 = \RecipeYield,
}

\NewDocumentCommand \RecipeNameAndYield {m}{% !!! Need a way to show "starred" (favorite) recipes??
	\pgfkeys{/RecipeNameAndYield, default, #1}%	
	% Put the recipe name in the Index, unless the user sets NoIdxName = 1:
	% (How about any non-empty instead of = 1)?
	\ifnum\NoIndexName=1 \relax\else\index{\RecipeName}\fi
	% Can add up to three other Index entries:
	\ifx\IdxA\empty\relax\else\index{\IdxA}\fi
	\ifx\IdxB\empty\relax\else\index{\IdxB}\fi
	\ifx\IdxC\empty\relax\else\index{\IdxC}\fi
	% Did the user ask us to set up a label for cross-reference?:
	\ifx\XRefLbl\empty\relax\else\RecipeLabel{\XRefLbl}\fi
	% Add a bookmark (only adds bookmark if hyperref is active):
	\ifnum\NoIndexName=1 \RecipeBookmark{\IdxA}\else\RecipeBookmark{\RecipeName}\fi
	\begin{center}%
		{\FontRecipeName{\RecipeName}}\par%	
		\ifx\RecipeYield\empty\relax\else {\textit{\RecipeYield}}\par \fi%	
	\end{center}%
}

%************ SECTION 5B - RECIPE STORY ****************************************************
\def\RecipeStoryIndent{20 pt}
\newbox\StoryBox
\newbox\StoryBoxA
\NewDocumentCommand \RecipeStory {o +m}  % arg1 = optional = lines for first part
{                                        % arg2 = story text
	\setbox\StoryBox\vbox
	{
		\FontRecipeStory
		\leftskip=\RecipeStoryIndent \rightskip=\leftskip	% Q 66332 see also Q 156986 (egreg and his comment)
		{#2\par}           % the \par needed by \lettrine
	}
	\IfNoValueTF {#1}    % NoValue=TRUE or NoValue=FALSE?
	{\unvbox\StoryBox} % NoValue=TRUE, so flush the full StoryBox
	{                  % handle the NoValue=FALSE case
		\ifnum #1 > 0      % > 0 means we split the story into two part; print the first part here
				{                % = 0 means ALL of story is deferred
					\setbox\StoryBoxA=\vsplit\StoryBox to #1\dimexpr \RStoryFontBaseline % 
					\unvbox\StoryBoxA  % 
				}
		\fi
	}
}

\NewDocumentCommand \FinishRecipeStory{}%
{\ifvoid\StoryBox \else \smallskip\unvbox\StoryBox\par\medskip\fi} %

%************ SECTION 5C - INGREDIENTS AND STEPS ENVIRONMENT *******************************
\usepackage{multicol}	% allows multiple columns in selected areas of a single-column document
\usepackage[strict]{changepage}	% for \adjustwidth	
\usepackage{xfp}				% for fpeval floating point macro

\pgfkeys{
	/IngredientsAndSteps/.is family, /IngredientsAndSteps,
	default/.style = {AdjIFont = 1, AdjIBaseline = 0,
			AdjSFont = 0, AdjSBaseline = 0,
			AddWidth = 0, RaggedCols   = 0},
	AdjIFont/.estore in     = \AdjustIFont,
	AdjIBaseline/.estore in = \AdjustIBaseline,
	AdjSFont/.estore in     = \AdjustSFont,
	AdjSBaseline/.estore in = \AdjustSBaseline,
	AddWidth/.estore in     = \AddPageWidth,
	RaggedCols/.estore in   = \RaggedColumns,
}

\newdimen\IFont  \newdimen\IBase  \newdimen\SFont  \newdimen\SBase  \newdimen\AddWidth

\NewDocumentEnvironment{IngredientsAndSteps}{o}		% !!! Need \raggedcolumns ??
% Do the before-environment setup:
{
	\IfNoValueTF {#1}%
	% Handle the simple case = no optional key/value argument:  
	{\AddWidth = 0 pt}
	% Handle the special case = there IS an optional key/value argument:
	{
		\pgfkeys{/IngredientsAndSteps, default, #1}%

		\IFont    = \AdjustIFont pt
		\IBase    = \AdjustIBaseline pt
		\SFont    = \AdjustSFont pt
		\SBase    = \AdjustSBaseline pt
		\AddWidth = \AddPageWidth pt

		\ifdim\IBase = 0 pt \IBase = \IFont \fi
		\ifdim\SFont = 0 pt \SFont = \IFont \fi
		\ifdim\SBase = 0 pt \SBase = \IFont \fi

		\RenewDocumentCommand \FontIngredients {}{\FontIngDefault\fontsize{\fpeval{\StdIFontSize * \IFont}}
			{\fpeval{\StdIBaseline * \IBase}}\selectfont}
		\RenewDocumentCommand \FontSteps {}{\FontStepsDefault\fontsize{\fpeval{\StdSFontSize * \SFont}}
			{\fpeval{\StdSBaseline * \SBase}}\selectfont}
	}
	% Now, finish setup:	
	\adjustwidth{-\AddWidth}{-\AddWidth}
	\begin{multicols}{2}%
		} % End of the before-environment setup
		% Now do the after-environment cleanup:
		{
	\end{multicols}
	\endadjustwidth
	% We need to restore the fonts to default size IF there was an optional key/value argument:
	\IfValueT {#1}
	{
		\RenewDocumentCommand \FontIngredients {}{\FontIngDefault}
		\RenewDocumentCommand \FontSteps       {}{\FontStepsDefault}
	}
} % End of the "after environment" cleanup

%************ SECTION 5D - COMMANDS TO DISPLAY INGREDIENTS AND STEPS ***********************
\NewDocumentCommand \IngredientsHeading {O{0} m O{2}}%
{\begingroup
	\setlength{\parindent}{0pt}
	\ifnum #1 > 0 {\vspace{#1 pt}}\fi
	\FontIngTitle\color{clrIngTitle} #2\par%
	\vspace{#3 pt}
	\endgroup}%

\NewDocumentCommand \ListIngredientsAndSteps {o +m +m}
{
	\IfValueT {#1}{\IngredientsHeading{#1}[3]}%
	{\begingroup\ifx\relax#2\relax\else\FontIngredients{}\IngredientsList{#2}\fi\endgroup}%
	{\begingroup\ifx\relax#3\relax\else\FontSteps{}\RecipeSteps{#3}\fi\endgroup}%
}

\NewDocumentCommand \IngredientsList { >{\SplitList{\par}} +m}{%
	\setlength{\parskip}{2pt}\raggedright%
	\ProcessList{#1}{\ProcessIngList}
}
\newcommand\ProcessIngList[1]{\hangindent1em #1\par}

\NewDocumentCommand \RecipeSteps { >{\SplitList{\par}} +m}{%         
	\begin{enumerate}[itemsep=2pt plus 1 pt minus 1pt, parsep=0pt plus 1pt,
			topsep=4.5pt plus 2.0pt minus 1.0pt, leftmargin=*]
		\ProcessList{#1}{\item}
	\end{enumerate}
}

%************ SECTION 5E - THE TIPS, CHEFFY AND CHEFNOTES ENVIRONMENTS *********************
\usepackage{tcolorbox}
\tcbset{fonttitle=\FontTitleColorBox}

\NewDocumentEnvironment{Tip}{+m}	%
{
	\begin{tcolorbox}[colback=clrBackTip,colframe=clrFrameTip, title=Food for Thought] %
		{#1}%
	\end{tcolorbox}
}

\NewDocumentEnvironment{Cheffy}{+m} %
{
	\begin{tcolorbox}[colback=clrBackCheffy,colframe=clrFrameCheffy, title=Let's Get Cheffy] %
		{#1}%
	\end{tcolorbox}
}

\NewDocumentEnvironment{ChefNotes}{+m}	%
{
	\begin{tcolorbox}[colback=clrBackNotes,colframe=clrFrameNotes,title=Chefnotes] %
		\ChefNote{}{#1}%
	\end{tcolorbox}
}

%************ SECTION 6 - ANY TIKZ DRAWINGS ************************************************
\usepackage{tikz}
\usetikzlibrary{calc}

\def\hatscale{0.05}
\def\hatwidth{12.5pt}
\NewDocumentCommand \MakeChefHat{}
{
	\begin{tikzpicture}[scale=\hatscale]
		\coordinate (HA) at (1.27, 0.635);
		\coordinate (HB) at (5.87375, 0);
		\coordinate (HC) at (1.5825, 0.9475);
		\coordinate (HD) at (5.77375, 0.15);
		\coordinate (HE) at (1.5875,1.905);
		\coordinate (HF) at (1.8,1.74625);
		\coordinate (HG) at (6.0325, 1.27);
		\coordinate (HH) at (6.22, 1.33);
		\coordinate (HI) at (6.45, 2.18);
		\coordinate (HJ) at (6.35, 2.70);
		\coordinate (HK) at (8.5725, 3.97);
		\coordinate (HL) at (6.03, 4.7625);
		\coordinate (HM) at (6.03, 3.651);
		\coordinate (HN) at (2.2225, 5.3975);
		\coordinate (HO) at (0.159, 4.603);
		\coordinate (HP) at (1.587, 3.55);
		\coordinate (HQ) at (2.5, 4.7625);
		\coordinate (HR) at (2.778, 2.08);
		\coordinate (HS) at (2.7, 4.1275);
		\coordinate (HT) at (4.0, 1.98);
		\coordinate (HU) at (4.286, 3.4925);
		\coordinate (HV) at (5.3, 1.7);
		\coordinate (HW) at (5.55, 2.85);
		\coordinate (HX) at (4.286, 5.0);
		%
		\path[fill=black] (HA) to [bend left=20] (HB) to (HD) to [bend right=20] (HC) to (HA);
		\path[fill=black] (HA) to (HC) to (HF) to (HE) to (HA);
		\path[fill=black] (HE) to (HF) to [bend left=20] (HG) to (HH) to [bend right=20] (HE);
		\path[fill=black] (HB) to (HD) to (HG) to (HH) to (HB);
		\path[fill=black] (HH) to (HI) to (HJ) to (HG);
		\draw[line width=5*\hatscale] (HI) to [bend right=50] (HK);
		\draw[line width=5*\hatscale] (HK) to [bend right=50] (HL);
		\draw[line width=5*\hatscale] (HL) to [bend right=45] (HX);
		\draw[line width=5*\hatscale] ($(HX) + (0.4, -.4)$) to [bend right=45] (HN);
		\path[fill=black] (HM) to ($(HL) -(0.15, 0)$) to ($(HL) + (0.15, 0)$) to (HM);
		\draw[line width=5*\hatscale] (HN) to [bend right=50] (HO);
		\draw[line width=5*\hatscale] (HO) to [bend right=50] (HP);
		\draw[line width=5*\hatscale]  ($(HF) +(0, 0.15)$) to (HP);
		\path[fill=black] (HQ) to ($(HN) -(0.15, 0)$) to ($(HN) + (0.15, 0)$) to (HQ);
		\path[fill=black] (HS) to ($(HR) -(0.15, 0)$) to ($(HR) + (0.15, 0)$) to (HS);
		\path[fill=black] (HU) to ($(HT) -(0.15, 0)$) to ($(HT) + (0.15, 0)$) to (HU);
		\path[fill=black] (HW) to ($(HV) -(0.15, 0)$) to ($(HV) + (0.15, 0)$) to (HW);
	\end{tikzpicture}
}

\def\ovenscale{0.2666}
\def\ovenwidth{18pt}
\NewDocumentCommand \MakeOven{} % inspired by the oven in the tikzsymbols package
{
	\begin{tikzpicture}[scale=\ovenscale]
		\draw [line width=2*\ovenscale] (0,0) rectangle (2,1.5);
		\draw [line width=2*\ovenscale] (0.125,0.125) rectangle (1.875,1);
		\draw [line width=2*\ovenscale] (0.75,1.125) rectangle (1.25,1.375);
		\draw [line width=2*\ovenscale] (0.5,0.875) to (1.5,0.875);
		\draw [fill=black] (0.23, 1.25) circle [radius=0.1];
		\draw [fill=black] (0.52, 1.25) circle [radius=0.1];
		\draw [fill=black] (1.48, 1.25) circle [radius=0.1];
		\draw [fill=black] (1.77, 1.25) circle [radius=0.1];
	\end{tikzpicture}
}

\ifMakeCookbook
\else
	\NewDocumentCommand \MakeTitlePage{}
	{
		\begingroup
		%\centering
		\begin{tikzpicture}[x=1in,y=1in]
			\coordinate (TA) at (0, 0);
			\coordinate (TB) at (6.125, 8);  % we ignore bottom 0.25 of text area to prevent newpage
			\coordinate (TC) at (1, 0.25);
			\coordinate (TD) at (5.125, 7.5);
			\coordinate (TE) at (1, 7.5);
			\coordinate (TF) at (5.125, 0.25);
			\useasboundingbox (TA) rectangle (TB);
			\draw [line width=1] (TC) to (TD);
			\node[anchor=north west, align=left, font=\FontTitlepageTitle] at (TE) {The Miller Family Cooks};
			\node[left, font=\FontTitlepageAuthor] at (TF) {Michaeljon Miller and Family};
		\end{tikzpicture}
		\endgroup
	}
\fi

%************ SECTION 7 - FRONT MATTER HELPER COMMANDS *************************************

\ifMakeCookbook
	%%%%%%%%%%%%%%%%%%%%%%%%%%%%%%%%%%%%%%%%%%%%%%%%%%%%%%%%%%%%%%%%%%%%%%%%
	% This group for the MAKECOOKBOOK bundle
	%%%%%%%%%%%%%%%%%%%%%%%%%%%%%%%%%%%%%%%%%%%%%%%%%%%%%%%%%%%%%%%%%%%%%%%%
	\NewDocumentCommand \MyTitlePage{}
	{
		\noindent\vspace{144pt}
		\begin{center}
			{
				{\FontTitlepageTitle The Miller Family Cooks} \par
				\vspace*{\fill}
				{\FontTitlepageAuthor by Michaeljon Miller and Family}\par
			}
		\end{center}
		\vspace{40pt} \clearpage
	}

	\NewDocumentCommand \ColophonFonts{}
	{
		Serif font (primary document font):
		EB Garamond, \Copyright{} The EB Garamond Project Authors\par
		\quad (github.com/octaviopardo/EBGaramond12). Licensed under the SIL Open Font License, Version 1.1.\par
		Sans Serif font:
		Lato, \Copyright{} Łukasz Dziedzic. Licensed under the SIL Open Font License, Version 1.1.\par
		Script font: Italianno, \Copyright{} TypeSETit, LLC. Licensed under the SIL Open Font License, Version 1.1.\par
	}

	\NewDocumentCommand \CopyrightInfo{}
	{Copyright \Copyright{} 2025, Michaeljon Miller, all rights reserved}

	\NewDocumentCommand \MyDedication{}
	{\hrule\vspace{5pt} To my family. And possibly some other stuff. \vspace{5pt}\hrule}

	\NewDocumentCommand \DedicationVerso{}{}

\else
	%%%%%%%%%%%%%%%%%%%%%%%%%%%%%%%%%%%%%%%%%%%%%%%%%%%%%%%%%%%%%%%%%%%%%%%%
	% This group for our cookbook
	%%%%%%%%%%%%%%%%%%%%%%%%%%%%%%%%%%%%%%%%%%%%%%%%%%%%%%%%%%%%%%%%%%%%%%%%
	\NewDocumentCommand \MyTitlePage{} {\noindent \MakeTitlePage{} \clearpage}

	\NewDocumentCommand \ColophonFonts{}
	{
		Serif font (primary document font):
		Adobe Garamond Pro, \Copyright{} Adobe Systems Incorporated.\par
		Sans Serif font:
		Adobe Myriad Pro, \Copyright{} Adobe Systems Incorporated.\par
		Script font: Italianno, \Copyright{} TypeSETit, LLC. Licensed under the SIL Open Font License, Version 1.1.\par
	}

	\NewDocumentCommand \CopyrightInfo{}
	{Copyright \Copyright{} 2025, Michaeljon Miller, all rights reserved}

	\NewDocumentCommand \MyDedication{}
	{\hrule\vspace{5pt}
		To my family.  Now when you ask \Quote{how do you make...}, I just need to give you a page number.%
		\vspace{5pt}\hrule}

	\NewDocumentCommand \DedicationVerso{}
	{
		\begin{center}
			{\FontChapterTitle
				One cannot think well, \\
				love well, sleep well, \\
				if one has not dined well.\\
				\vspace{5pt}
				\hrule
				\vspace{5pt}
				Uno non pu\`o pensare bene, \\
				amare bene, dormire bene, \\
				se  non ha mangiato bene. \\
				\vspace{5pt}
				\hrule
				\vspace{5pt}
				Virginia Woolf
			}
		\end{center}
	}
\fi

%************ SECTION 8 - GRAPHICX RELATED COMMANDS ****************************************
\graphicspath{{images/}{../images/}} % look in the img directory (subdir of book root or sibling of tex)
\usepackage[skip=2pt, labelformat=empty, font={rm,it}]{caption} % <-
\captionsetup{justification=centering} % this is needed to have multi-line captions centered

\pgfkeys{
	/SideBySide/.is family, /SideBySide,
	default/.style = {VertAlign = c, LeftCaption = \empty, RightCaption = \empty},
	VertAlign/.estore in    = \VerticalAlign,
	LeftCaption/.estore in  = \LeftText,
	RightCaption/.estore in = \RightText,
}

\NewDocumentCommand \SideBySide{O{\empty} m m} % Q 5769
{
	\pgfkeys{/SideBySide, default, #1}%
	%
	\begin{figure}[htb]
		\centering
		\begin{minipage}[\VerticalAlign]{0.49\textwidth} % align at t= top, c = center, b = bottom
			\centering
			\includegraphics[width=0.97\textwidth]{#2} %
			\ifx\LeftText\empty\relax\else\caption{\LeftText}\fi
		\end{minipage}\hfill%
		\begin{minipage}[\VerticalAlign]{0.49\textwidth}
			\centering
			\includegraphics[width=0.97\textwidth]{#3} % 
			\ifx\RightText\empty\relax\else\caption{\RightText}\fi
		\end{minipage}
	\end{figure}
}

%************ SECTION 9 - HYPERREF RELATED COMMANDS ****************************************
\ifHyperRef
	\usepackage[bookmarks=true,colorlinks=true, allcolors=clrHyperRef]{hyperref}
	\usepackage{bookmark} % Q 247158 shows Heiko (the author) loads both hyperref and bookmark

	\newcounter{ctrRecipe} % initially set to zero by \newcounter macro
	\NewDocumentCommand \RecipeBookmark{m}{%
		\ifx\%1\empty\relax\else\stepcounter{ctrRecipe}\pdfbookmark[1]{#1}{Recipe\arabic{ctrRecipe}}\fi
	}
\else
	\providecommand\phantomsection{}% Q 44088
	\NewDocumentCommand \RecipeBookmark{m}{}%
\fi

\NewDocumentCommand \RecipeLabel{m}{\phantomsection\label{recipe:#1}}
\NewDocumentCommand \RecipePageNo{m}{\pageref{recipe:#1}}

\NewDocumentCommand \CookbookIndex{}
{
	\cleardoublepage % flush all material and clear until you start new odd numbered (recto) page
	%
	% the following line needed if hyperref is loaded (p. 18 of hyperref manual):
	\phantomsection\addcontentsline{toc}{chapter}{\indexname} % see also Q 59619
	%
	\printindex
}

